% Pandoc template for the fed_fraud documents.
% Based on proposal_gen/templates/detailed_proposal_template/detailed_1_classic.tex
% (Classic report): centered title page, table of contents, ruled section
% headings and the Geodesic house footer. Fonts are Latin Modern (lmodern).
% Two layouts: the full report (title page and contents) and, with
% -V brief, a compact one-page-brief header in the same style.
\documentclass[letterpaper,11pt]{article}
\usepackage[T1]{fontenc}
\usepackage[utf8]{inputenc}
\usepackage{lmodern}
\usepackage[letterpaper,top=0.9in,bottom=1.2in,left=1in,right=1in,footskip=0.6in]{geometry}
\usepackage[table]{xcolor}
\usepackage{geodesicbrand}
\usepackage{amsmath,amssymb}
\usepackage{enumitem}
\usepackage{array}
\usepackage{calc}
\usepackage{longtable}
\usepackage{makecell}
\usepackage[most]{tcolorbox}
\usepackage{titlesec}
\usepackage{needspace}
\usepackage{etoolbox}
\usepackage{hyperref}
\hypersetup{colorlinks=true,urlcolor=coolblack,linkcolor=coolblack,filecolor=geoblue,
  pdftitle={Forensic methods handbook},pdfauthor={\geofirm}}
\urlstyle{same}

\setlist{itemsep=2pt,leftmargin=1.4em}
\setlength{\parindent}{0pt}
\setlength{\parskip}{6pt}
\providecommand{\tightlist}{\setlength{\itemsep}{0pt}\setlength{\parskip}{0pt}}
\setcounter{secnumdepth}{2}

% Section style: house rule under each heading
\titleformat{\section}{\Large\bfseries\color{coolblack}}{\thesection}{0.6em}{}[{\color{coolblack}\titlerule[0.5pt]}]
\titleformat{\subsection}{\large\bfseries\color{geomid}}{\thesubsection}{0.6em}{}
\titleformat{\subsubsection}{\normalsize\bfseries\color{geomid}}{\thesubsubsection}{0.6em}{}
\titlespacing*{\section}{0pt}{16pt}{8pt}

% Tables are written by latex/filters.lua (navy header, banded rows)
\setlength{\LTpre}{6pt}\setlength{\LTpost}{6pt}

\begin{document}

% TITLE PAGE
\begin{titlepage}
\newcommand{\HRule}{\rule{\linewidth}{0.3mm}}
\begin{center}
\vspace*{0.2in}
\textsc{\LARGE \geofirm}\\[0.6cm]
\includegraphics[width=30mm]{icon.png}\\[8mm]
\textsc{\Large Federal Fraud Expansion\\[0.15cm] Working Document}\\[0.5cm]
\HRule\\[0.4cm]
{\huge\bfseries Forensic methods handbook}\\[0.3cm]
{\large The mathematics behind each fraud test}\\[0.2cm]
\HRule\\[1cm]
\begin{minipage}{0.8\textwidth}\centering\large
Authored by:\\[0.18cm]
\geosigner, \geosignertitle: \geoemail
\end{minipage}\\[0.8cm]
{\large When: \compiletime}\\[1.2cm]
\geosurface{catenoid}{4.6cm}\\[0.9cm]
$\mathcal{G}(x)=x^x\qquad\nabla_{\dot\gamma}\dot\gamma=0$\\[3mm]
{\small\itshape \geopromise}
\end{center}
\end{titlepage}

\pagestyle{geodesic}
\tableofcontents
\newpage

\textbf{Status:} Draft v0.1 (2026-11-03), for Steve's review. Companion
to the \href{methodology.pdf}{methodology}, which says how a lead is
scored, validated and referred; this handbook says exactly how each test
is computed, when it is valid and what it cannot show. Every formula
here has a tested reference implementation in
\texttt{analytics/forensic/} (run on synthetic data only; no real entity
appears in the tests).

\hypertarget{bottom-line}{%
\section*{Bottom line}\label{bottom-line}}
\addcontentsline{toc}{section}{Bottom line}

Fraud leaves arithmetic traces: digits that people invent, amounts
steered just under a limit, invoices paid twice, bids that cluster too
neatly, providers who bill more hours than a day has, and ``separate''
companies that share a phone line. This handbook turns each trace into a
written test with a formula, a threshold, a false-positive story and a
validation plan. Ten families of tests cover every fraud area in the
roadmap and in the \href{../universe/all_agencies.pdf}{all-agency plan}:

\begin{enumerate}
\def\labelenumi{\arabic{enumi}.}
\tightlist
\item
  \textbf{Digit and number tests} (Benford's law, round numbers,
  duplication).
\item
  \textbf{Duplicate, split and impossible-billing tests.}
\item
  \textbf{Peer outliers, growth and threshold bunching}, with Bayesian
  shrinkage so small providers are not over-flagged.
\item
  \textbf{Collusion screens} for sealed bids.
\item
  \textbf{Entity resolution} (Fellegi-Sunter record linkage).
\item
  \textbf{Network analysis} with hub discounting and risk propagation.
\item
  \textbf{Eligibility and identity tests} (death, incarceration,
  impossible identities).
\item
  \textbf{Improper payment estimation} with confidence intervals, so we
  can check and extend agency estimates.
\item
  \textbf{Reconciliation and audit-trail tests} across USAspending,
  Treasury, PaymentAccuracy.gov and Single Audit data.
\item
  \textbf{Score combination and validation} (likelihood ratios, lift,
  false discovery control).
\end{enumerate}

The rule that keeps this honest: \textbf{a test result is a reason to
look, not a finding.} Only an Inspector General or DOJ confirms fraud,
and only a confirmed dollar counts as savings.

\hypertarget{where-each-method-sits}{%
\section*{1. Where each method sits}\label{where-each-method-sits}}
\addcontentsline{toc}{section}{1. Where each method sits}

Tests work at five levels. Public data supports the first three in Phase
1; the last two mostly need internal records reached through OIGs in
Phase 2.

\begingroup\small\setlength{\tabcolsep}{4pt}\renewcommand{\arraystretch}{1.25}
\rowcolors{2}{white}{geolight}
\begin{longtable}[l]{@{}>{\raggedright\arraybackslash}p{\dimexpr 0.097\linewidth-2\tabcolsep}>{\raggedright\arraybackslash}p{\dimexpr 0.204\linewidth-2\tabcolsep}>{\raggedright\arraybackslash}p{\dimexpr 0.265\linewidth-2\tabcolsep}>{\raggedright\arraybackslash}p{\dimexpr 0.265\linewidth-2\tabcolsep}>{\raggedright\arraybackslash}p{\dimexpr 0.168\linewidth-2\tabcolsep}@{}}
\rowcolor{coolblack}\leavevmode\color{white}\bfseries Level & \leavevmode\color{white}\bfseries Unit tested & \leavevmode\color{white}\bfseries Example tests & \leavevmode\color{white}\bfseries Phase 1 public data & \leavevmode\color{white}\bfseries Phase 2 internal data \\
\endhead
Population & All awards or payments of an agency or program & Benford, round numbers, bunching, reconciliation & USAspending award and account files, PaymentAccuracy.gov, Treasury MTS & General ledger, payment files \\
Entity & One contractor, grantee, provider, retailer & Peer outliers, growth, new-entity-large-award, entity-level Benford & USAspending, SAM.gov, CMS files, SNAP retailer data & Claims, invoices \\
Network & Groups of entities & Shared attributes, components, risk propagation & SAM.gov addresses and contacts, NPPES, Form 990, state registries & Bank accounts, devices, IP addresses \\
Transaction & One payment, claim, invoice or voucher & Duplicates, splits, impossible days, after-death payments & Rarely public & Payment, claim, card and voucher files \\
Estimate & The program's improper payment rate & Stratified and ratio estimators, monetary unit sampling & Published agency estimates & Sample case files \\
\end{longtable}\endgroup

\hypertarget{ground-rules-for-forensic-testing}{%
\section*{2. Ground rules for forensic
testing}\label{ground-rules-for-forensic-testing}}
\addcontentsline{toc}{section}{2. Ground rules for forensic testing}

\begin{itemize}
\tightlist
\item
  \textbf{Hypothesis first.} Each test is written down with the scheme
  it targets before it is run. A test without a scheme is data dredging.
\item
  \textbf{Pre-registration.} Thresholds, peer groups and exclusions are
  fixed in the screen's specification before results are seen, and
  changes are logged with reasons.
\item
  \textbf{Multiple testing.} Running one test per entity across 100,000
  entities guarantees thousands of chance ``hits''. Entity-level
  p-values are controlled with the Benjamini-Hochberg procedure (section
  12.4), not a fixed 0.05.
\item
  \textbf{Practical significance.} With millions of records, a
  chi-square test rejects for trivially small deviations. Every test
  reports an effect size (MAD, ratio, excess dollars) alongside its
  p-value, and effect size decides whether anyone looks.
\item
  \textbf{Innocent explanations first.} Each test lists its known benign
  causes and checks the main ones before a result enters a lead package.
\item
  \textbf{Reproducibility.} Every result cites its data pull (source,
  date, hash) and the code version, so another analyst can rerun it
  exactly. This follows the GAO Yellow Book's evidence standard and the
  AICPA Statement on Standards for Forensic Services No.~1.
\item
  \textbf{Known error rates.} Each screen's back-test (section 12.3)
  gives its measured hit rate and false-positive rate. That is what an
  expert needs to defend a method in court (the Daubert standard asks
  for a known or potential error rate).
\end{itemize}

\hypertarget{digit-and-number-tests}{%
\section*{3. Digit and number tests}\label{digit-and-number-tests}}
\addcontentsline{toc}{section}{3. Digit and number tests}

\hypertarget{benfords-law}{%
\subsection*{3.1 Benford's law}\label{benfords-law}}
\addcontentsline{toc}{subsection}{3.1 Benford's law}

Many naturally occurring sets of amounts (invoices, claims, award
values) follow Benford's law: the leading digit \(d\) appears with
probability

\[P(D_1 = d) = \log_{10}\left(1 + \frac{1}{d}\right), \quad d = 1, \dots, 9,\]

so about 30.1\% of amounts start with 1 and only 4.6\% start with 9. The
first-two-digits version, with \(d = 10, \dots, 99\), uses the same
formula and is the sharper forensic test. The second digit has
probability

\[P(D_2 = d) = \sum_{k=1}^{9} \log_{10}\left(1 + \frac{1}{10k + d}\right), \quad d = 0, \dots, 9.\]

\textbf{Why it detects fraud.} People inventing numbers spread leading
digits too evenly and favor middle digits; schemes that keep amounts
under a limit create spikes at specific first-two digits (for example 48
and 49 below a \$5,000 approval level).

\textbf{When it applies.} Benford's law holds for amounts that span
several orders of magnitude, arise from multiplying many factors, and
have no built-in floor or ceiling. It does \textbf{not} apply to
assigned numbers (invoice or account numbers), prices set on a schedule
(fixed per-diem rates, Medicare fee schedules), amounts with caps (a
\$10,000 grant maximum), or small samples. A thousand or more records
works best for first-two-digit tests.

\textbf{Statistics.} With \(n\) amounts, observed digit shares
\(\hat p_d\) and expected shares \(p_d\):

\begin{itemize}
\tightlist
\item
  Chi-square: \(\chi^2 = n \sum_d (\hat p_d - p_d)^2 / p_d\) with 8
  (first digit) or 89 (first two) degrees of freedom. Sensitive to
  \(n\); use for small populations.
\item
  Per-digit z with continuity correction:
  \(z_d = \left(|\hat p_d - p_d| - \tfrac{1}{2n}\right) / \sqrt{p_d(1-p_d)/n}\).
  Identifies which digits spike.
\item
  Mean absolute deviation:
  \(\text{MAD} = \frac{1}{k}\sum_d |\hat p_d - p_d|\), which does not
  grow with \(n\) and is the main conformity measure.
\end{itemize}

\begingroup\small\setlength{\tabcolsep}{4pt}\renewcommand{\arraystretch}{1.25}
\rowcolors{2}{white}{geolight}
\begin{longtable}[l]{@{}>{\raggedright\arraybackslash}p{\dimexpr 0.163\linewidth-2\tabcolsep}>{\raggedright\arraybackslash}p{\dimexpr 0.204\linewidth-2\tabcolsep}>{\raggedright\arraybackslash}p{\dimexpr 0.204\linewidth-2\tabcolsep}>{\raggedright\arraybackslash}p{\dimexpr 0.163\linewidth-2\tabcolsep}>{\raggedright\arraybackslash}p{\dimexpr 0.265\linewidth-2\tabcolsep}@{}}
\rowcolor{coolblack}\leavevmode\color{white}\bfseries Test & \leavevmode\color{white}\bfseries Close conformity & \leavevmode\color{white}\bfseries Acceptable & \leavevmode\color{white}\bfseries Marginal & \leavevmode\color{white}\bfseries Nonconformity \\
\endhead
First digit & MAD up to 0.006 & 0.006 to 0.012 & 0.012 to 0.015 & Above 0.015 \\
Second digit & Up to 0.008 & 0.008 to 0.010 & 0.010 to 0.012 & Above 0.012 \\
First two digits & Up to 0.0012 & 0.0012 to 0.0018 & 0.0018 to 0.0022 & Above 0.0022 \\
\end{longtable}\endgroup

Thresholds are Nigrini's (2012), derived from large real-world data
sets.

\textbf{How we use it.} At the population level, nonconformity is a
reason to look at the program's controls, not at any one payee. At the
entity level (one vendor's 400 invoices, one provider's claims), each
entity's MAD is compared with its peers' distribution of MADs, and only
the extreme tail with enough records is reviewed. Benford is a
\textbf{weak signal alone}; it never counts as one of the two required
signals unless combined with a digit spike at a known threshold.

\hypertarget{companion-number-tests}{%
\subsection*{3.2 Companion number tests}\label{companion-number-tests}}
\addcontentsline{toc}{subsection}{3.2 Companion number tests}

\begingroup\small\setlength{\tabcolsep}{4pt}\renewcommand{\arraystretch}{1.25}
\rowcolors{2}{white}{geolight}
\begin{longtable}[l]{@{}>{\raggedright\arraybackslash}p{\dimexpr 0.168\linewidth-2\tabcolsep}>{\raggedright\arraybackslash}p{\dimexpr 0.335\linewidth-2\tabcolsep}>{\raggedright\arraybackslash}p{\dimexpr 0.268\linewidth-2\tabcolsep}>{\raggedright\arraybackslash}p{\dimexpr 0.229\linewidth-2\tabcolsep}@{}}
\rowcolor{coolblack}\leavevmode\color{white}\bfseries Test & \leavevmode\color{white}\bfseries Formula & \leavevmode\color{white}\bfseries What it catches & \leavevmode\color{white}\bfseries Benign causes \\
\endhead
Last-two-digits & Chi-square of the last two digits against uniform (1\% each) & Invented cents, or amounts all ending .00 or .99 & Whole-dollar pricing, fixed fees \\
Round numbers & Share of amounts that are exact multiples of \$1,000 (or \$100), tested against a clean reference share with a one-sided binomial test & Estimates billed as actuals; made-up invoices & Lump-sum milestone payments, fixed grants \\
Number duplication & Most frequent exact amounts, ranked & One odd amount (say \$4,987.13) paid 60 times & Recurring contracts, subscriptions \\
Summation & Share of total dollars by first-two digits; each of the 90 groups should hold about 1/90 & A few very large amounts driving the total & One large legitimate contract \\
Relative size factor & Largest payment divided by second largest, per vendor or cardholder & One inflated or misplaced payment & Annual true-ups \\
\end{longtable}\endgroup

\hypertarget{duplicate-split-and-impossible-billing-tests}{%
\section*{4. Duplicate, split and impossible-billing
tests}\label{duplicate-split-and-impossible-billing-tests}}
\addcontentsline{toc}{section}{4. Duplicate, split and
impossible-billing tests}

These are transaction tests. They need payment, invoice, card or voucher
records, so most run in Phase 2, but the award-level versions run in
Phase 1 on USAspending contract actions.

\hypertarget{duplicate-payments}{%
\subsection*{4.1 Duplicate payments}\label{duplicate-payments}}
\addcontentsline{toc}{subsection}{4.1 Duplicate payments}

Duplicate payments are the most common recoverable error and a common
fraud (resubmitting an invoice with a changed number). The test runs in
four tiers, strongest first:

\begin{enumerate}
\def\labelenumi{\arabic{enumi}.}
\tightlist
\item
  \textbf{Exact:} same vendor, amount, invoice number and date.
\item
  \textbf{Same invoice, different date:} same vendor, amount and
  \emph{normalized} invoice number. Normalization uppercases, strips
  prefixes (INV, \#) and suffixes (-A, -R, DUP), maps look-alike
  characters (O to 0, I and l to 1, S to 5) and removes leading zeros.
\item
  \textbf{Near-identical invoice:} same vendor and amount, normalized
  invoice numbers within one edit (Damerau-Levenshtein distance 1, which
  counts a transposed pair of digits as one edit), paid within 90 days.
\item
  \textbf{Same invoice, different vendor:} same amount and invoice
  number paid to two vendor records, which points to duplicate vendor
  master entries or a diverted payment.
\end{enumerate}

Expected false positives are recurring identical charges (monthly rent,
subscriptions); these are listed per vendor and excluded after review.

\hypertarget{split-purchases-and-split-awards}{%
\subsection*{4.2 Split purchases and split
awards}\label{split-purchases-and-split-awards}}
\addcontentsline{toc}{subsection}{4.2 Split purchases and split awards}

Rules that require more scrutiny above a dollar limit (the FAR
micro-purchase threshold for card purchases, the simplified acquisition
threshold for contracts, internal approval levels) invite splitting. The
test groups purchases by the same buyer from the same vendor within a
short window (three days for cards, 30 days for contract actions) and
flags any group where \textbf{every} item is under the limit but the
\textbf{sum} is over it.

At the population level, the same idea becomes a \textbf{bunching test}
(section 5.4): an excess of amounts just under the limit.

The limits change with inflation adjustments (the October 2025 FAR
adjustment raised the micro-purchase threshold to \$15,000 and the
simplified acquisition threshold to \$350,000); each screen reads the
limit in force on the transaction date from FAR 2.101 rather than a
fixed number.

\hypertarget{impossible-days-and-impossible-volumes}{%
\subsection*{4.3 Impossible days and impossible
volumes}\label{impossible-days-and-impossible-volumes}}
\addcontentsline{toc}{subsection}{4.3 Impossible days and impossible
volumes}

For any biller paid by time or by service (Criminal Justice Act panel
attorneys, workers' compensation providers, Medicaid personal care
aides, time-and-materials contractors), total billed hours per person
per day across all matters is

\[H_{b,t} = \sum_{m} h_{b,m,t}.\]

\(H > 24\) is impossible; reviews also look at \(H > 16\) on many days
and at service counts that exceed what one clinician can deliver (for
example, more 60-minute psychotherapy sessions in a day than hours in a
working day). The test is simple, explains itself to a jury, and has
produced many health care fraud prosecutions.

\hypertarget{services-and-payments-after-death}{%
\subsection*{4.4 Services and payments after
death}\label{services-and-payments-after-death}}
\addcontentsline{toc}{subsection}{4.4 Services and payments after death}

Any service or payment dated after the beneficiary's date of death,
beyond a short grace period for legitimate end-of-month payments, is
flagged. Death data come from SSA through Treasury's Do Not Pay service
(which began receiving SSA's full death file in 2023 under a three-year
pilot) or from state vital records; the Phase 1 public alternative is
the obituary-backed case library, which is too sparse to screen with.

\hypertarget{peer-outliers-growth-and-threshold-bunching}{%
\section*{5. Peer outliers, growth and threshold
bunching}\label{peer-outliers-growth-and-threshold-bunching}}
\addcontentsline{toc}{section}{5. Peer outliers, growth and threshold
bunching}

\hypertarget{peer-groups}{%
\subsection*{5.1 Peer groups}\label{peer-groups}}
\addcontentsline{toc}{subsection}{5.1 Peer groups}

Every outlier test compares an entity with \textbf{peers in the same
program, service, place and period}: a hospice with hospices in its
county and year, a SNAP retailer with stores of the same type within a
mile. A national comparison mostly finds specialty, case mix and size.

\hypertarget{robust-z-scores}{%
\subsection*{5.2 Robust z-scores}\label{robust-z-scores}}
\addcontentsline{toc}{subsection}{5.2 Robust z-scores}

For a peer group with values \(x_i\),

\[z_i = \frac{x_i - \operatorname{median}(x)}{1.4826 \cdot \operatorname{MAD}(x)},\]

where
\(\operatorname{MAD}(x) = \operatorname{median}|x_i - \operatorname{median}(x)|\).
The constant makes the scale match a standard deviation for normal data,
and unlike the mean and standard deviation, a few extreme entities
cannot mask themselves by inflating the scale. Skewed measures (dollars,
counts) are log-transformed first. Review threshold: \(z > 3.5\) (set
per screen in the calibration week).

\hypertarget{count-data-and-empirical-bayes-shrinkage}{%
\subsection*{5.3 Count data and empirical Bayes
shrinkage}\label{count-data-and-empirical-bayes-shrinkage}}
\addcontentsline{toc}{subsection}{5.3 Count data and empirical Bayes
shrinkage}

For counts (claims, services, meals) with exposure (beneficiaries,
enrolled children), let entity \(i\) have observed count \(y_i\) and
expected count \(E_i\) from its peers' rate. The raw ratio \(y_i / E_i\)
is unstable when \(E_i\) is small: 3 claims against 0.5 expected looks
like a 6x outlier but is weak evidence. We model

\[y_i \sim \text{Poisson}(\theta_i E_i), \qquad \theta_i \sim \text{Gamma}(\alpha, \beta),\]

fit \(\alpha, \beta\) across the peer population by the method of
moments, and rank on the posterior

\[\theta_i \mid y_i \sim \text{Gamma}(\alpha + y_i,\ \beta + E_i), \qquad
\mathbb{E}[\theta_i \mid y_i] = \frac{\alpha + y_i}{\beta + E_i}.\]

Leads are ranked by the posterior probability that \(\theta_i\) exceeds
a material multiple of peers (1.5x by default), which balances size of
excess against strength of evidence. This is the same approach used in
disease mapping and in the hierarchical models behind CMS's hospital
readmission and mortality measures.

\hypertarget{bunching-below-thresholds}{%
\subsection*{5.4 Bunching below
thresholds}\label{bunching-below-thresholds}}
\addcontentsline{toc}{subsection}{5.4 Bunching below thresholds}

If a limit \(T\) invites steering, amounts pile up just below it. Bin
amounts in width \(w\), fit a polynomial of degree 5 to bin counts
outside a window around \(T\), and predict counts inside the window.
Excess mass is

\[B = \sum_{\text{bins below } T \text{ in window}} (c_j - \hat c_j), \qquad
b = B \Big/ \sum \hat c_j,\]

following Chetty et al.~(2011). A ratio \(b\) well above zero, with
matching missing mass just above \(T\), is strong population-level
evidence of splitting or steering. Applied to card purchases
(micro-purchase limit), contract actions (simplified acquisition and
certified cost data thresholds) and grant draws (prior-approval
thresholds). Standard errors come from bootstrapping the residuals.

\hypertarget{rapid-growth-and-change-points}{%
\subsection*{5.5 Rapid growth and change
points}\label{rapid-growth-and-change-points}}
\addcontentsline{toc}{subsection}{5.5 Rapid growth and change points}

For an entity's monthly payments \(x_t\) with a peer-based expected
level \(\mu_0\), the one-sided CUSUM

\[S_t = \max(0,\ S_{t-1} + x_t - \mu_0 - k)\]

signals when \(S_t > h\). Setting \(k\) to half the shift worth
detecting and \(h\) at four to five standard deviations gives a low
false-alarm rate. CUSUM catches sudden ramps (a new hospice billing at
the cap within months, a retailer whose redemptions triple) that annual
totals hide.

\hypertarget{concentration}{%
\subsection*{5.6 Concentration}\label{concentration}}
\addcontentsline{toc}{subsection}{5.6 Concentration}

The Gini coefficient of an entity's billing across service codes, or of
a program's payments across entities, measures concentration. A provider
whose billing collapses onto one or two high-paying codes (Gini near 1
where peers sit near 0.5) is an upcoding signal; a program where a
handful of new entities take most of the money is a capture signal.

\hypertarget{collusion-screens-for-sealed-bids}{%
\section*{6. Collusion screens for sealed
bids}\label{collusion-screens-for-sealed-bids}}
\addcontentsline{toc}{section}{6. Collusion screens for sealed bids}

Bid rigging leaves statistical traces in the full set of bids, which
state DOTs publish as bid tabulations (DOT-1) and some agencies release
on request. For each tender with bids
\(b_{(1)} \le b_{(2)} \le \dots \le b_{(n)}\):

\begingroup\small\setlength{\tabcolsep}{4pt}\renewcommand{\arraystretch}{1.25}
\rowcolors{2}{white}{geolight}
\begin{longtable}[l]{@{}>{\raggedright\arraybackslash}p{\dimexpr 0.171\linewidth-2\tabcolsep}>{\raggedright\arraybackslash}p{\dimexpr 0.400\linewidth-2\tabcolsep}>{\raggedright\arraybackslash}p{\dimexpr 0.429\linewidth-2\tabcolsep}@{}}
\rowcolor{coolblack}\leavevmode\color{white}\bfseries Screen & \leavevmode\color{white}\bfseries Formula & \leavevmode\color{white}\bfseries Collusive pattern \\
\endhead
Coefficient of variation & \(\text{CV} = \sigma_b / \bar b\) & Low: cover bids cluster near an agreed price \\
Spread & \((b_{(n)} - b_{(1)}) / b_{(1)}\) & Low \\
Relative distance & \(\text{RD} = (b_{(2)} - b_{(1)}) / \sigma_{\text{losing}}\) & High: a gap between the designated winner and tight cover bids \\
Skewness and kurtosis & Third and fourth standardized moments of the bids & Distorted by clustered cover bids \\
Cover ratio & Mean losing bid divided by winning bid & Stable across tenders for the same group \\
\end{longtable}\endgroup

Imhof, Karagök and Rutz (2018) found these distribution screens separate
known Swiss cartels from competitive markets, and Huber and Imhof (2019)
raised accuracy by combining them in a classifier. Our use: compute the
screens per tender, then compare each regular group of bidders with the
same work type and region.

\textbf{Rotation test.} Among a stable group of firms, cartels often
take turns. Count how often the same firm wins two consecutive tenders,
and compare with a permutation distribution formed by shuffling the
order of the same wins. Too \emph{few} repeats (an orderly rotation) is
the signal.

\textbf{Market allocation.} A win-share matrix shows, for each pair of
firms that often bid together, who wins. Firms that meet often but each
win only in ``their'' county or season are flagged for review.

\textbf{Benign causes.} Few bidders in rural markets, standard
engineer's estimates that anchor bids, capacity constraints. Screens are
leads for DOJ's Procurement Collusion Strike Force, not findings.

\hypertarget{entity-resolution}{%
\section*{7. Entity resolution}\label{entity-resolution}}
\addcontentsline{toc}{section}{7. Entity resolution}

The entity spine joins records that describe the same business or person
across SAM.gov, USAspending, NPPES, Form 990, state registries and the
watch lists. Errors here create false networks (merging strangers) or
hide real ones (splitting one entity).

\textbf{Blocking.} Compare only pairs that share a cheap key (ZIP plus
name token, phone, email domain, UEI, NPI), which reduces billions of
pairs to millions.

\textbf{Comparison.} For each candidate pair, compare fields to get an
agreement vector \(\gamma\): name (Jaro-Winkler similarity at least 0.92
after removing legal suffixes), street address (USPS-standardized,
exact), phone (exact), officer name, email domain.

\textbf{Fellegi-Sunter weights.} With
\(m_k = P(\gamma_k = 1 \mid \text{match})\) and
\(u_k = P(\gamma_k = 1 \mid \text{non-match})\), the match weight is

\[W = \sum_k \left[\gamma_k \log_2 \frac{m_k}{u_k} + (1 - \gamma_k)\log_2\frac{1 - m_k}{1 - u_k}\right].\]

A rare field agreement (same phone number) earns a large weight because
\(u_k\) is tiny; a common one (same city) earns little. \(m\) and \(u\)
are estimated without labels by the EM algorithm and checked on the 500
hand-labeled pairs in the methodology.

\textbf{Decision.} Pairs above an upper threshold are matches; below a
lower threshold, non-matches; between them, clerical review. The upper
threshold is the lowest weight with labeled precision of 95\% or better.

\textbf{Clustering.} Linking matched pairs transitively can chain
unrelated records (A matches B, B matches C, A and C are strangers). We
cluster on the weighted match graph and split clusters whose weakest
internal link falls below the upper threshold.

\hypertarget{network-analysis}{%
\section*{8. Network analysis}\label{network-analysis}}
\addcontentsline{toc}{section}{8. Network analysis}

\hypertarget{shared-attribute-graph}{%
\subsection*{8.1 Shared-attribute graph}\label{shared-attribute-graph}}
\addcontentsline{toc}{subsection}{8.1 Shared-attribute graph}

Entities connect through shared attributes: address, phone, email,
officer, registered agent and, in Phase 2, bank account and device. We
build the bipartite entity-attribute graph and project it onto entities
with an Adamic-Adar style weight:

\[w_{ij} = \sum_{a \in A_i \cap A_j} \frac{1}{\log(1 + \deg(a))}.\]

A phone shared by two firms counts far more than an address shared by
4,000 firms at a registered agent. Attributes on the benign-hub list, or
shared by more entities than a set cap, are dropped.

\hypertarget{co-location-test}{%
\subsection*{8.2 Co-location test}\label{co-location-test}}
\addcontentsline{toc}{subsection}{8.2 Co-location test}

Is an address hosting more entities than its type explains? If addresses
of that type (single-family home, small office, mailbox store) average
\(\lambda\) registered entities, the Poisson tail

\[P(X \ge k) = 1 - \sum_{i=0}^{k-1} \frac{e^{-\lambda}\lambda^i}{i!}\]

gives how surprising \(k\) entities are. Twelve federal contractors at a
house where 1.5 are typical is a lead; twelve at a coworking space is
not.

\hypertarget{components-brokers-and-communities}{%
\subsection*{8.3 Components, brokers and
communities}\label{components-brokers-and-communities}}
\addcontentsline{toc}{subsection}{8.3 Components, brokers and
communities}

\begin{itemize}
\tightlist
\item
  \textbf{Connected components} above an edge-weight floor give
  candidate rings.
\item
  \textbf{Betweenness centrality} finds brokers (one consultant,
  preparer or billing agent linking many outliers), a pattern in E-Rate,
  SBA lending and health care cases.
\item
  \textbf{Community detection} (Louvain or Leiden) splits large
  components into tighter groups for review.
\item
  \textbf{Temporal motifs} catch waves: many entities registered within
  weeks of each other at linked addresses, just before a program opens.
\end{itemize}

\hypertarget{risk-propagation}{%
\subsection*{8.4 Risk propagation}\label{risk-propagation}}
\addcontentsline{toc}{subsection}{8.4 Risk propagation}

Personalized PageRank spreads risk from known bad actors (excluded,
debarred, convicted entities) through weighted links. With restart
probability \(\alpha\) and seed distribution \(s\),

\[r = \alpha s + (1 - \alpha) P^{\top} r,\]

where \(P\) is the row-normalized weighted adjacency matrix. An entity's
score measures how close it sits to known fraud through strong links.
Used as a signal, never as guilt by association: a high score must be
paired with an independent signal of its own.

\hypertarget{eligibility-and-identity-tests}{%
\section*{9. Eligibility and identity
tests}\label{eligibility-and-identity-tests}}
\addcontentsline{toc}{section}{9. Eligibility and identity tests}

Mostly Phase 2, run on applicant and payee files under OIG data
agreements.

\begingroup\small\setlength{\tabcolsep}{4pt}\renewcommand{\arraystretch}{1.25}
\rowcolors{2}{white}{geolight}
\begin{longtable}[l]{@{}>{\raggedright\arraybackslash}p{\dimexpr 0.178\linewidth-2\tabcolsep}>{\raggedright\arraybackslash}p{\dimexpr 0.411\linewidth-2\tabcolsep}>{\raggedright\arraybackslash}p{\dimexpr 0.411\linewidth-2\tabcolsep}@{}}
\rowcolor{coolblack}\leavevmode\color{white}\bfseries Test & \leavevmode\color{white}\bfseries Rule & \leavevmode\color{white}\bfseries Programs \\
\endhead
Death match & Payment or service date after date of death (SSA full death file, state records) & Medicare, Medicaid, FEHB, VA, PBGC, RRB, SNAP, OPM annuities \\
Incarceration match & Benefit paid during a period of incarceration (BOP and state prison data) & SSI, UI, SNAP, VA \\
Impossible identity & SSN issued after the person's stated birth year plus a plausible gap, or SSN issued after the claim; SSNs issued after June 2011 are randomized and need SSA verification instead of area-number rules & UI, SBA, student aid \\
Identity reuse & One SSN, bank account, phone, email or device behind claims in several states or programs & UI, SNAP, Medicaid, SBA, ED \\
Synthetic identity & Thin credit-file-style profile: new email, prepaid phone, no history, shared device & UI, SBA, tax credits \\
Concurrent eligibility & Benefits that cannot legally be received together (for example some disability and earnings combinations) & SSA, RRB, DOL, VA \\
\end{longtable}\endgroup

Bank-account reuse is the strongest single identity signal in pandemic
cases; it is the core of cross-agency idea X-2 in the roadmap.

\hypertarget{estimating-improper-payments}{%
\section*{10. Estimating improper
payments}\label{estimating-improper-payments}}
\addcontentsline{toc}{section}{10. Estimating improper payments}

Agencies estimate improper payments by sampling payments and reviewing
them (Payment Integrity Information Act of 2019; OMB Circular A-123,
Appendix C). We need this math for three jobs: reading published
estimates critically, designing samples for Phase 2 pilots, and
estimating the dollars behind a screen's leads without overstating them.

\hypertarget{stratified-expansion-estimator}{%
\subsection*{10.1 Stratified expansion
estimator}\label{stratified-expansion-estimator}}
\addcontentsline{toc}{subsection}{10.1 Stratified expansion estimator}

Payments are split into strata \(h\) (by size, region or payment type).
From \(N_h\) payments, \(n_h\) are sampled and reviewed, giving improper
dollars \(y_{hi}\). The estimated total and its variance are

\[\hat Y = \sum_h N_h \bar y_h, \qquad
\widehat{\operatorname{Var}}(\hat Y) = \sum_h N_h^2 \left(1 - \frac{n_h}{N_h}\right)\frac{s_h^2}{n_h}.\]

The improper payment rate is \(\hat R = \hat Y / X\), where \(X\) is
total outlays, with standard error \(\operatorname{SE}(\hat Y)/X\). A
two-sided interval at confidence \(1-a\) is
\(\hat Y \pm z_{1-a/2}\operatorname{SE}\).

\hypertarget{ratio-estimator}{%
\subsection*{10.2 Ratio estimator}\label{ratio-estimator}}
\addcontentsline{toc}{subsection}{10.2 Ratio estimator}

When improper dollars scale with payment size, the combined ratio
estimator is more precise:

\[\hat R = \frac{\sum_h N_h \bar y_h}{\sum_h N_h \bar x_h}, \qquad
\widehat{\operatorname{Var}}(\hat R) = \frac{1}{\hat X^2} \sum_h N_h^2 \left(1 - \frac{n_h}{N_h}\right)\frac{s_{d,h}^2}{n_h},\]

where \(x\) is dollars paid, \(\hat X = \sum_h N_h \bar x_h\) and
\(d = y - \hat R x\) are the residuals (Taylor linearization).

\hypertarget{when-the-normal-interval-fails}{%
\subsection*{10.3 When the normal interval
fails}\label{when-the-normal-interval-fails}}
\addcontentsline{toc}{subsection}{10.3 When the normal interval fails}

Improper dollars are highly skewed: most sampled payments are fully
correct, a few are wholly wrong. Normal-theory intervals then cover the
truth less often than stated. In our simulation (8\% of payments
improper, lognormal payment sizes; see
\texttt{analytics/tests/test\_forensic.py}), nominal 90\% intervals
covered the true rate about 77\% of the time at \(n = 100\), 84\% at
\(n = 400\) and 88\% at \(n = 1{,}500\). Remedies, in order of
preference: larger samples in high-variance strata (Neyman allocation,
\(n_h \propto N_h S_h\)), bootstrap intervals, and for audits with few
errors, monetary unit sampling.

\hypertarget{monetary-unit-sampling}{%
\subsection*{10.4 Monetary unit sampling}\label{monetary-unit-sampling}}
\addcontentsline{toc}{subsection}{10.4 Monetary unit sampling}

Monetary unit (dollar unit) sampling selects dollars, not payments, so
large payments are more likely to be sampled. With book value \(BV\),
\(n\) sampled dollars and taintings \(t_{(1)} \ge t_{(2)} \ge \dots\)
(each error's share of its payment's value), the Stringer upper error
limit is

\[\text{UEL} = \frac{BV}{n}\left[U(0) + \sum_{i \ge 1} \big(U(i) - U(i-1)\big)\, t_{(i)}\right],\]

where \(U(k) = \tfrac{1}{2}\chi^2_{c}(2k + 2)\) is the upper Poisson
confidence limit after \(k\) errors (\(U(0) = 3.00\) and \(U(1) = 4.74\)
at 95\%). The sample size that keeps the bound under a tolerable
misstatement \(TM\) when no errors are expected is
\(n = BV \cdot U(0) / TM\).

\hypertarget{conventions-to-know}{%
\subsection*{10.5 Conventions to know}\label{conventions-to-know}}
\addcontentsline{toc}{subsection}{10.5 Conventions to know}

\begin{itemize}
\tightlist
\item
  HHS-OIG and Medicare contractors extrapolate overpayments using the
  lower limit of a confidence interval (HHS-OIG uses the lower limit of
  a two-sided 90\% interval), which favors the provider. Our lead
  packages use the same convention when they size dollars.
\item
  HHS-OIG's RAT-STATS is the standard tool for these samples; our code
  reproduces its core estimators for checking.
\item
  \textbf{Unknown payments} (payments that cannot be shown proper for
  lack of documentation) are reported separately from improper payments
  under A-123 and are neither fraud nor savings.
\item
  \textbf{Improper is not fraud.} GAO estimated federal fraud losses at
  \$233 billion to \$521 billion a year (GAO-24-105833) using a
  different method; FY2025 improper payment estimates totaled about
  \$186 billion across 64 programs at 15 agencies (GAO-26-108694).
  Neither figure is ever used as a savings claim.
\end{itemize}

\hypertarget{checking-a-published-estimate}{%
\subsection*{10.6 Checking a published
estimate}\label{checking-a-published-estimate}}
\addcontentsline{toc}{subsection}{10.6 Checking a published estimate}

From public reports we can test whether an estimate is internally
consistent: whether its stated interval matches its sample size and
strata, whether the rate moved more than its standard error would allow
from one year to the next without a stated cause, and whether programs
that ``did not report'' left large susceptible outlays unmeasured. These
checks become questions in a Phase 2 proposal, not findings.

\hypertarget{reconciliation-and-audit-trail-tests}{%
\section*{11. Reconciliation and audit-trail
tests}\label{reconciliation-and-audit-trail-tests}}
\addcontentsline{toc}{section}{11. Reconciliation and audit-trail tests}

\hypertarget{public-reconciliations-phase-1}{%
\subsection*{11.1 Public reconciliations (Phase
1)}\label{public-reconciliations-phase-1}}
\addcontentsline{toc}{subsection}{11.1 Public reconciliations (Phase 1)}

\begingroup\small\setlength{\tabcolsep}{4pt}\renewcommand{\arraystretch}{1.25}
\rowcolors{2}{white}{geolight}
\begin{longtable}[l]{@{}>{\raggedright\arraybackslash}p{\dimexpr 0.109\linewidth-2\tabcolsep}>{\raggedright\arraybackslash}p{\dimexpr 0.297\linewidth-2\tabcolsep}>{\raggedright\arraybackslash}p{\dimexpr 0.297\linewidth-2\tabcolsep}>{\raggedright\arraybackslash}p{\dimexpr 0.297\linewidth-2\tabcolsep}@{}}
\rowcolor{coolblack}\leavevmode\color{white}\bfseries Check & \leavevmode\color{white}\bfseries Left side & \leavevmode\color{white}\bfseries Right side & \leavevmode\color{white}\bfseries What a gap means \\
\endhead
Account to award & USAspending File A/B obligations and outlays by Treasury account & File C award-level obligations linked to D1/D2 awards & Unlinked or unreported awards; spending invisible to payee-level screens \\
Agency totals & USAspending agency outlays & Treasury Monthly Treasury Statement outlays & Reporting scope differences (trust funds, offsetting receipts); explains, rarely fraud \\
Program outlays & PaymentAccuracy.gov susceptible outlays by program & USAspending outlays by Assistance Listing & Programs measured on a different base than they pay out \\
Prime to subaward & Prime award obligation & Sum of FFATA subawards reported under it & Subawards exceeding the prime are impossible and point to reporting errors or double-counting; missing subawards hide pass-throughs \\
Grantee self-report & Grant drawdowns in USAspending & Federal expenditures reported in the grantee's Single Audit & Large unexplained gaps between what was drawn and what was audited \\
Single Audit findings & Federal Audit Clearinghouse findings and questioned costs & The grantee's continuing awards & Repeat material weaknesses or questioned costs with no change in funding \\
\end{longtable}\endgroup

The Single Audit data (fac.gov) are the broadest public internal-control
data in government: every non-federal entity spending \$1 million or
more of federal awards a year (the threshold since the 2024 Uniform
Guidance revision) files one. Two cheap signals cover every grant-making
agency: repeat findings at the same grantee, and \textbf{auditor
concentration} (one small audit firm auditing many grantees that share
officers, or whose audits never report findings).

\hypertarget{audit-trail-tests-on-internal-ledgers-phase-2}{%
\subsection*{11.2 Audit-trail tests on internal ledgers (Phase
2)}\label{audit-trail-tests-on-internal-ledgers-phase-2}}
\addcontentsline{toc}{subsection}{11.2 Audit-trail tests on internal
ledgers (Phase 2)}

These follow the journal-entry and fraud procedures in auditing
standards (AU-C 240 and PCAOB AS 2401) and are run on the general
ledger, vendor master, payment and card files that OIGs can obtain.

\begingroup\small\setlength{\tabcolsep}{4pt}\renewcommand{\arraystretch}{1.25}
\rowcolors{2}{white}{geolight}
\begin{longtable}[l]{@{}>{\raggedright\arraybackslash}p{\dimexpr 0.241\linewidth-2\tabcolsep}>{\raggedright\arraybackslash}p{\dimexpr 0.380\linewidth-2\tabcolsep}>{\raggedright\arraybackslash}p{\dimexpr 0.380\linewidth-2\tabcolsep}@{}}
\rowcolor{coolblack}\leavevmode\color{white}\bfseries Test & \leavevmode\color{white}\bfseries Rule & \leavevmode\color{white}\bfseries Scheme \\
\endhead
Sequence gaps and duplicates & Missing or repeated check, invoice and obligation numbers & Removed or fabricated documents \\
Three-way match & Payment without a matching order and receiving report, or with amounts that differ & Payment for goods not delivered \\
Vendor-employee match & Vendor address, phone or bank account equal to an employee's & Shell vendors owned by insiders \\
Bank account change & Vendor bank account changed shortly before a large payment, especially after an email request & Business email compromise, diversion \\
Segregation of duties & Same user created the vendor, entered the invoice and approved the payment & Insider fraud \\
Timing & Entries on weekends, holidays or after hours; entries just after period close & Concealment, manipulation \\
Just-below-approval & Entries clustered under a user's approval limit (bunching test, section 5.4) & Splitting to avoid review \\
Card tests & Split purchases (section 4.2), merchant category codes outside policy, purchases on leave days & Purchase and travel card misuse (Government Charge Card Abuse Prevention Act of 2012 requires OIG risk assessments of these programs) \\
Payroll & Employees with no leave ever, shared bank accounts, pay after separation & Ghost employees \\
\end{longtable}\endgroup

\hypertarget{combining-signals-calibrating-and-validating}{%
\section*{12. Combining signals, calibrating and
validating}\label{combining-signals-calibrating-and-validating}}
\addcontentsline{toc}{section}{12. Combining signals, calibrating and
validating}

\hypertarget{likelihood-ratios}{%
\subsection*{12.1 Likelihood ratios}\label{likelihood-ratios}}
\addcontentsline{toc}{subsection}{12.1 Likelihood ratios}

For a binary signal \(S\), the case library gives

\[\text{LR}^{+} = \frac{P(S \mid \text{known case})}{P(S \mid \text{population})},\]

estimated with additive smoothing so an unobserved combination never
gets a ratio of zero. Under conditional independence, posterior odds are
prior odds times the product of the signals' ratios:

\[\frac{P(F \mid S_1, \dots, S_k)}{1 - P(F \mid S_1, \dots, S_k)} = \frac{\pi}{1 - \pi}\prod_{j} \text{LR}_j.\]

Correlated signals (shared address and shared phone) are grouped and
only the strongest in each group enters the product.

\hypertarget{the-base-rate-trap}{%
\subsection*{12.2 The base-rate trap}\label{the-base-rate-trap}}
\addcontentsline{toc}{subsection}{12.2 The base-rate trap}

Even a strong signal leaves most flagged entities innocent when fraud is
rare. A signal present in 40\% of known cases and 0.5\% of the
population has \(\text{LR}^{+} = 80\); with a 1-in-1,000 prior, the
posterior is only about 7\%. This is why the methodology requires two
independent signals and human review, and why scores are presented as
rankings for review, not as probabilities of guilt. (The numbers in this
paragraph are illustrative and are reproduced in the scoring tests.)

\hypertarget{validation-metrics}{%
\subsection*{12.3 Validation metrics}\label{validation-metrics}}
\addcontentsline{toc}{subsection}{12.3 Validation metrics}

\begin{itemize}
\tightlist
\item
  \textbf{Lift at the top \(q\):} share of known cases in the top \(q\)
  of scores divided by their share overall. Release bar: lift of 5 or
  more at the top 5\% (methodology section 5).
\item
  \textbf{Precision at \(k\):} share of known cases (or of hand-review
  ``worth a look'' judgments) among the top \(k\) leads, matched to
  investigator capacity.
\item
  \textbf{Average precision:} area under the precision-recall curve;
  preferred over ROC AUC when cases are rare.
\item
  \textbf{Calibration:} when scores are shown as probabilities in Phase
  2, isotonic regression on held-out cases maps scores to observed
  rates.
\item
  \textbf{Temporal holdout:} score the population as it looked before
  each case became public, so a screen is not credited for finding cases
  it learned from.
\end{itemize}

\hypertarget{controlling-false-discoveries}{%
\subsection*{12.4 Controlling false
discoveries}\label{controlling-false-discoveries}}
\addcontentsline{toc}{subsection}{12.4 Controlling false discoveries}

When the same test runs on \(m\) entities, rank their p-values
\(p_{(1)} \le \dots \le p_{(m)}\) and flag the largest \(k\) with

\[p_{(k)} \le \frac{k}{m}\, q.\]

The Benjamini-Hochberg procedure keeps the expected share of chance hits
among flagged entities at or below \(q\) (we use \(q = 0.05\) for
entity-level digit and outlier tests).

\hypertarget{methods-by-fraud-area}{%
\section*{13. Methods by fraud area}\label{methods-by-fraud-area}}
\addcontentsline{toc}{section}{13. Methods by fraud area}

Each fraud area in the roadmap maps to a small set of tests. ``P1''
means the test runs on public data in Phase 1; ``P2'' needs internal
data.

\begingroup\small\setlength{\tabcolsep}{4pt}\renewcommand{\arraystretch}{1.25}
\rowcolors{2}{white}{geolight}
\begin{longtable}[l]{@{}>{\raggedright\arraybackslash}p{\dimexpr 0.127\linewidth-2\tabcolsep}>{\raggedright\arraybackslash}p{\dimexpr 0.218\linewidth-2\tabcolsep}>{\raggedright\arraybackslash}p{\dimexpr 0.218\linewidth-2\tabcolsep}>{\raggedright\arraybackslash}p{\dimexpr 0.218\linewidth-2\tabcolsep}>{\raggedright\arraybackslash}p{\dimexpr 0.218\linewidth-2\tabcolsep}@{}}
\rowcolor{coolblack}\leavevmode\color{white}\bfseries Fraud area & \leavevmode\color{white}\bfseries Main schemes & \leavevmode\color{white}\bfseries Primary tests (section) & \leavevmode\color{white}\bfseries Data & \leavevmode\color{white}\bfseries Agencies (examples) \\
\endhead
Health care billing & Upcoding, phantom services, kickback-driven referrals, hospice and home health mills & Peer outliers with shrinkage (5.3), concentration (5.6), growth (5.5), impossible days (4.3, P2), after-death services (4.4, P2), networks (8) & CMS provider and utilization files (P1); claims (P2) & HHS, VA, OPM (FEHB), DOL (FECA), USPS (workers' comp) \\
Procurement fronts & Pass-through set-aside fronts, shell bidders & Entity resolution (7), shared attributes and co-location (8.1-8.2), new entity with large award, subaward reconciliation (11.1) & SAM.gov, FPDS, USAspending subawards & DoD, VA, DOT, GSA, SBA, USACE, all agencies (U4) \\
Bid rigging & Rotation, cover bids, market allocation & Collusion screens and rotation test (6) & State DOT bid tabs; agency bid data (P2) & DOT, USACE, DoD, GSA \\
Price and cost mischarging & Inflated prices, defective pricing, labor mischarging & Unit price outliers (5.2), modification growth (U3), bunching under cost-data thresholds (5.4), impossible hours (4.3, P2) & FPDS, USAspending; invoices and timesheets (P2) & DoD, NASA, DOE, GSA \\
Grant misuse & Duplicate funding, unallowable costs, ghost participants & Text-similarity duplicate awards, Single Audit findings and auditor concentration (11.1), cost per participant outliers (5.2), drawdown gaps (11.1) & USAspending, fac.gov, Form 990, SBIR.gov & NSF, NIH, DOE, EPA, NEA, NEH, IMLS, regional commissions, LSC, CPB \\
Loans and guarantees & Identity and income fraud, lender and agent rings, straw borrowers & Entity resolution (7), broker centrality (8.3), digit tests on loan amounts (3), bunching under program caps (5.4) & SBA and USDA loan-level releases, EXIM authorizations, HMDA & SBA, USDA, HUD (FHA), EXIM, DFC \\
Benefits & Ineligible, deceased or incarcerated recipients; identity reuse & Death, incarceration and identity tests (9), sampling estimates (10) & Aggregate public statistics (P1); payee files (P2) & SSA, RRB, DOL, VA, OPM, USDA \\
Tax credits & Preparer-driven false claims & Geographic outliers with shrinkage (5.3) on IRS statistics by ZIP; preparer networks (P2) & IRS SOI ZIP data (P1); returns (P2, section 6103 limits) & Treasury \\
Customs and trade & Undervaluation, false origin, transshipment & Mirror-trade gaps, importer entity resolution, co-location (8.2) & Census trade data, manifests & DHS \\
Subsidies (E-Rate, Lifeline, BEAD) & Consultant-steered bids, phantom subscribers & Collusion screens (6), broker centrality (8.3), duplicate subscribers (4.1, P2) & USAC open data & FCC, Commerce \\
Disaster and emergency aid & Duplicate claims, ineligible addresses & Duplicate tests (4.1), identity reuse (9), address validation & FEMA public data (P1); applicant files (P2) & DHS (FEMA), SBA, HUD (CDBG-DR) \\
Settlement and refund distributions & Synthetic claimants, duplicate claims & Duplicate and identity tests (4.1, 9) on claims files & Distribution reports (P1); claims (P2) & CFPB, FTC, SEC, CFTC, DOJ \\
Internal (every agency) & Card misuse, payroll ghosts, insider vendors, payment diversion & Split purchases (4.2), audit-trail tests (11.2), bunching (5.4) & OIG charge card risk assessments (P1); ledgers and card files (P2) & All bodies (U7) \\
\end{longtable}\endgroup

\hypertarget{reference-implementation}{%
\section*{14. Reference implementation}\label{reference-implementation}}
\addcontentsline{toc}{section}{14. Reference implementation}

\texttt{analytics/forensic/} holds a small, readable implementation of
every test in this handbook, with 30 unit tests that run on synthetic
data:

\begingroup\small\setlength{\tabcolsep}{4pt}\renewcommand{\arraystretch}{1.25}
\rowcolors{2}{white}{geolight}
\begin{longtable}[l]{@{}>{\raggedright\arraybackslash}p{\dimexpr 0.302\linewidth-2\tabcolsep}>{\raggedright\arraybackslash}p{\dimexpr 0.698\linewidth-2\tabcolsep}@{}}
\rowcolor{coolblack}\leavevmode\color{white}\bfseries Module & \leavevmode\color{white}\bfseries Sections \\
\endhead
\texttt{digits.py} & 3: Benford first, second, first-two and last-two digit tests; MAD conformity; per-digit z; summation; round numbers; duplication; relative size factor \\
\texttt{duplicates.py} & 4: invoice normalization, four-tier duplicate test, split purchases, impossible days, after-death services \\
\texttt{outliers.py} & 5: robust z, Poisson-Gamma shrinkage, CUSUM, bunching, Gini \\
\texttt{bids.py} & 6: tender screens, rotation test, win-share matrix \\
\texttt{linkage.py} & 7: Fellegi-Sunter EM, thresholds, Jaro-Winkler \\
\texttt{network.py} & 8: shared-attribute projection, components, personalized PageRank, co-location \\
\texttt{sampling.py} & 10: stratified and ratio estimators, Neyman allocation, sample sizes, monetary unit sampling \\
\texttt{scoring.py} & 12: likelihood ratios, posterior odds, lift, precision, average precision, two-signal rule \\
\end{longtable}\endgroup

Run
\texttt{python3\ -m\ unittest\ discover\ -s\ analytics/tests\ -t\ analytics}
from the repository root. Production pipelines may use faster libraries
but must reproduce these functions' results on the test fixtures.

\hypertarget{sources}{%
\section*{Sources}\label{sources}}
\addcontentsline{toc}{section}{Sources}

\begin{itemize}
\tightlist
\item
  Nigrini, M. (2012). \emph{Benford's Law: Applications for Forensic
  Accounting, Auditing, and Fraud Detection.} Wiley.
\item
  Fellegi, I. and Sunter, A. (1969). ``A Theory for Record Linkage.''
  \emph{Journal of the American Statistical Association} 64(328).
\item
  Winkler, W. (1988). ``Using the EM Algorithm for Weight Computation in
  the Fellegi-Sunter Model of Record Linkage.'' Proceedings of the
  Section on Survey Research Methods, American Statistical Association.
\item
  Imhof, D., Karagök, Y. and Rutz, S. (2018). ``Screening for Bid
  Rigging: Does It Work?'' \emph{Journal of Competition Law \&
  Economics} 14(2).
\item
  Huber, M. and Imhof, D. (2019). ``Machine Learning with Screens for
  Detecting Bid-Rigging Cartels.'' \emph{International Journal of
  Industrial Organization} 65.
\item
  Chetty, R., Friedman, J., Olsen, T. and Pistaferri, L. (2011).
  ``Adjustment Costs, Firm Responses, and Micro vs.~Macro Labor Supply
  Elasticities.'' \emph{Quarterly Journal of Economics} 126(2).
\item
  Benjamini, Y. and Hochberg, Y. (1995). ``Controlling the False
  Discovery Rate.'' \emph{Journal of the Royal Statistical Society B}
  57(1).
\item
  Cochran, W. (1977). \emph{Sampling Techniques}, 3rd ed.~Wiley.
\item
  \href{https://www.gao.gov/products/gao-15-593sp}{GAO-15-593SP: A
  Framework for Managing Fraud Risks in Federal Programs}
\item
  \href{https://www.gao.gov/products/GAO-24-105833}{GAO-24-105833:
  federal fraud losses estimated at \$233B to \$521B a year}
\item
  \href{https://www.gao.gov/products/gao-26-108694}{GAO-26-108694:
  FY2025 improper payments}
\item
  \href{https://paymentaccuracy.gov/agencies-and-programs}{PaymentAccuracy.gov
  agencies and programs}
\item
  \href{https://www.whitehouse.gov/wp-content/uploads/2021/03/M-21-19.pdf}{OMB
  M-21-19: Appendix C to Circular A-123}
\item
  \href{https://www.fac.gov/}{Federal Audit Clearinghouse}
\item
  \href{https://fiscal.treasury.gov/data-transparency/DAIMS-current.html}{USAspending
  DATA Act file definitions (DAIMS)}
\item
  \href{https://oig.hhs.gov/compliance/rat-stats/}{HHS-OIG RAT-STATS}
\item
  \href{https://www.aicpa-cima.com/resources/download/statement-on-standards-for-forensic-services}{AICPA
  Statement on Standards for Forensic Services No.~1}
\item
  \href{https://www.justice.gov/procurement-collusion-strike-force}{DOJ
  Procurement Collusion Strike Force}
\end{itemize}

\end{document}
